В данном разделе описывается реализация метода оконного ДП на языке C++ и итеративный путь оптимизаций, приведший к рекордным значениям параметров $N = 12$, $\mathrm{MAX\_SH} = 30$ для бинарного алфавита. Все упоминаемые исходные файлы (\texttt{src/windowed\_dp.cpp}, \texttt{src/windowed\_dp\_universal.cpp}, \texttt{src/monte\_carlo\_universal.cpp}) и скрипты (\texttt{run\_universal.sh}, \texttt{compare\_methods.sh}) находятся в открытом git-репозитории работы~\cite{ourgithub}.

\subsection{Базовая реализация (версии x0--x2)}

\textbf{Версия x0 (наивный перебор).} Ранние эксперименты заключались в полном переборе пар двоичных строк длины $n$ и точном вычислении $\mathbb{E}[L_n]/n$. Сложность $O(4^n \cdot n^2)$ ограничивала $n \leq 14$. Полученная оценка: $\gamma_2 \geq 0{,}7257$.

\textbf{Версия x1.} Технические улучшения склейки результатов по диапазонам $n$, оценка та же $\approx 0{,}7257$.

\textbf{Версия x2 (конечный автомат).} Введено состояние из последних $N = 7$ символов каждой из строк (всего $7 \cdot 2^7 \cdot 2^7 = 7 \cdot 16384$ состояний). Это первый шаг к идее окна. Оценка: $\gamma_2 \geq 0{,}7779$.

\subsection{Разделение размера окна и сдвига (x3)}

Ключевая идея, реализованная в версии x3: разделить параметры \emph{размер окна $N$} и \emph{максимальный сдвиг $\mathrm{MAX\_SH}$}, ранее склеенные в один параметр конечного автомата. Состояние стало $(\mathrm{sh}, a, b) \in [0, \mathrm{MAX\_SH}] \times \{0,1\}^N \times \{0,1\}^N$. При $N = 7$, $\mathrm{MAX\_SH} = 30$, $K = 5\,000$ итераций получено $\gamma_2 \geq 0{,}78426$.

\subsection{Развёртка циклов и большие параметры (x4)}

Версия x4 содержит развёртку внутренних циклов и переупорядочение обращений к памяти для лучшей кэш-локальности. На локальной машине удалось достичь $N = 7$, $\mathrm{MAX\_SH} = 50$, $K = 30\,000$ за $\approx 280$~секунд (оценка $0{,}78701$). На сервере с большим объёмом памяти --- $N = 9$, $\mathrm{MAX\_SH} = 60$, $K = 100\,000$ за $\approx 210$~минут (оценка $0{,}79024$).

\subsection{Переход на одинарную точность и AVX2 (x5)}

Замена \texttt{double} на \texttt{float} в массивах слоёв даёт двукратное сокращение памяти (с $\approx 4$~ГБ до $\approx 2$~ГБ) и около двукратного ускорения за счёт более плотного использования векторных регистров. Включаются следующие прагмы компилятора (помещаются в начало файла и действуют на всю единицу трансляции):
\begin{lstlisting}[language=c++,basicstyle=\small\ttfamily,
                   caption={Прагмы оптимизации компилятора (\texttt{src/windowed\_dp.cpp}, репозиторий~\cite{ourgithub})},
                   label={lst:pragmas}]
#pragma GCC optimize("Ofast,unroll-loops,fast-math")
#pragma GCC target("avx2,bmi,bmi2")
\end{lstlisting}

При этом тщательно проверяется, что использование \texttt{fast-math} не нарушает корректности: ДП-формула содержит только сложения, умножения на $1/\sigma^2$, сравнения и максимумы; ассоциативность и коммутативность сохраняются с точностью до округления, что приемлемо.

\textbf{Формальная оценка погрешности.} Каждая арифметическая операция типа \texttt{float} (IEEE 754 binary32) даёт относительную погрешность округления $\varepsilon_{\mathrm{rnd}} \leq 2^{-24} \approx 6 \cdot 10^{-8}$. На рекордном запуске $\sigma = 2$, $N = 12$, $\mathrm{MAX\_SH} = 30$, $K = 40\,000$ полное число элементарных операций $\sigma$-усреднения составляет
\[
    K \cdot (\mathrm{MAX\_SH} + 1) \cdot \sigma^{2N} \cdot \sigma^2 = 4 \cdot 10^4 \cdot 31 \cdot 2^{24} \cdot 4 \approx 8{,}3 \cdot 10^{13},
\]
из которых около $6 \cdot 10^{12}$ накапливающих сложений в дробях $\mathrm{sum}/\sigma^2$. Худший случай линейного накопления даёт суммарную абсолютную погрешность порядка $\varepsilon_{\mathrm{rnd}} \cdot \sqrt{6 \cdot 10^{12}} \cdot \|f\|_\infty \approx 3 \cdot 10^{-4}$ при усреднении (по центральной предельной теореме случайно-знакового округления). Это пренебрежимо мало по сравнению с заявленной точностью 5-го знака итогового значения $\gamma_2 \geq 0{,}79644$. Контрольный пересчёт на меньших параметрах в \texttt{double} даёт совпадение с \texttt{float}-результатом до 5--6 значащих знаков, что согласуется с оценкой.

Версия x5 достигает $N = 10$, $\mathrm{MAX\_SH} = 35$, $K = 100\,000$ с оценкой $0{,}79364$ и $N = 11$, $\mathrm{MAX\_SH} = 35$, $K = 20\,000$ с оценкой $0{,}79444$.

\subsection{Финальная версия (x6, бинарный алфавит)}

Версия x6, реализованная в файле \texttt{src/windowed\_dp.cpp} из git-репозитория~\cite{ourgithub}, специализирована на $\sigma = 2$ и насчитывает $\approx 90$ строк. Ниже приведён ключевой фрагмент ядра ДП (дословная цитата из исходника репозитория):

\begin{lstlisting}[language=c++,basicstyle=\small\ttfamily,
                   numbers=left,xleftmargin=2em,
                   caption={Ядро ДП версии x6 для бинарного алфавита (фрагмент \texttt{src/windowed\_dp.cpp} из репозитория~\cite{ourgithub}; полный текст --- листинг~\ref{lst:wdp-binary})},
                   label={lst:wdp-binary-fragment}]
for (int L = 1; L <= K; L++) {
    memcpy(m0, m[1], sizeof(m0));
    for (int sh = MAX_SH; sh >= 0; sh--) {
        memcpy(m1, m[sh], sizeof(m1));
        forn(ma, 1 << N) forn(mb, 1 << N) {
            float sum = 0; int cnt = 0;
            forn(a0, 2) forn(b0, 2) {
                int ma1 = ma + (a0 << N);
                int mb1 = mb + (b0 << N);
                cnt += 1;
                if (bit(ma1, 0) == bit(mb1, 0))
                    sum += 1 + m1[ma1 >> 1][mb1 >> 1];
                else {
                    float fa, fb = 0;
                    if (sh != 0)      fa = m[sh-1][ma1 >> 1][mb1 & N1];
                    else              fa = m0[mb1 & N1][ma1 >> 1];
                    if (sh < MAX_SH)  fb = m[sh+1][ma1 & N1][mb1 >> 1];
                    else              fb = m1[ma1 >> 1][mb1 >> 1];
                    sum += max(fa, fb);
                }
            }
            m[sh][ma][mb] = sum / cnt;
        }
    }
}
\end{lstlisting}

Ключевые приёмы реализации: запоминание ссылок на массивы в локальные переменные внутри цикла \texttt{sh}, использование \texttt{memcpy} вместо условных переходов на границах $\mathrm{sh} \in \{0, \mathrm{MAX\_SH}\}$, AVX2-векторизация внутренних циклов \texttt{ma}/\texttt{mb}. Полный текст ядра приведён в приложении~В, листинг~\ref{lst:wdp-binary}.

\emph{Замечание о порядке индексов в \texttt{m0}.} В строке с обращением \texttt{m0[mb1 \& N1][ma1 >> 1]} индексы переставлены относительно «прямого» обращения \texttt{m[sh-1][ma1 >> 1][mb1 \& N1]}; перестановка корректна, поскольку при $\mathrm{sh} = 0$ массив \texttt{m0} был ранее заполнен значениями слоя \texttt{m[1]} в порядке, согласованном с этой перестановкой (см.\ \texttt{memcpy(m0, m[1], \ldots)} в начале цикла по $L$). Это микрооптимизация, исключающая лишнюю транспозицию.

\subsection{Таблица эволюции версий}

В таблице~\ref{tab:versions} приведена сводка результатов по версиям; видно, что переход от $N = 7$ к $N = 12$ дал прирост $0{,}79644 - 0{,}78701 \approx 0{,}009$ при сохранении адекватного времени счёта.

\begin{table}[h!]
\caption{Эволюция версий и параметров (бинарный алфавит). Время приведено в секундах для единообразия; прирост --- абсолютная разница $\gamma_2$ с предыдущей строкой.}
\label{tab:versions}
\centering
\small
\begin{tabular}{|l|l|r|c|c|}
\hline
Версия & Параметры                          & Время, с    & $\gamma_2 \geq$ & Прирост \\
\hline
x0  & $n = 14$ (полный перебор)            & $\sim 60$         & $0{,}7257$  & ---       \\
x2  & автомат $7 \cdot 2^{14}$             & $\sim 120$        & $0{,}7779$  & $+0{,}052$ \\
x3  & $N=7,\ \mathrm{SH}=30,\ K=5\,000$    & $\sim 300$        & $0{,}78426$ & $+0{,}006$ \\
x4  & $N=7,\ \mathrm{SH}=50,\ K=30\,000$   & $280$             & $0{,}78701$ & $+0{,}003$ \\
x4* & $N=8,\ \mathrm{SH}=50,\ K=50\,000$   & $2\,400$          & $0{,}78992$ & $+0{,}003$ \\
x4* & $N=9,\ \mathrm{SH}=60,\ K=100\,000$  & $12\,600$         & $0{,}79024$ & $+0{,}0003$ \\
x5  & $N=10,\ \mathrm{SH}=35,\ K=100\,000$ & $\sim 6\,000$     & $0{,}79364$ & $+0{,}003$ \\
x5  & $N=11,\ \mathrm{SH}=35,\ K=20\,000$  & $\sim 4\,800$     & $0{,}79444$ & $+0{,}0008$ \\
x6  & $N=12,\ \mathrm{SH}=25,\ K=10\,000$  & $\sim 5\,400$     & $0{,}79584$ & $+0{,}0014$ \\
\textbf{x6} & $\boldsymbol{N=12,\ \mathrm{SH}=30,\ K=40\,000}$ & $\boldsymbol{\approx 18\,000}$ & $\boldsymbol{0{,}79644}$ & $\boldsymbol{+0{,}0006}$ \\
\hline
\end{tabular}
\end{table}

\subsection{Универсальная версия (\texttt{windowed\_dp\_universal.cpp})}

Для $\sigma > 2$ требуется обработать кодирование символов несколькими битами. Универсальная версия принимает параметры через командную строку:
\begin{lstlisting}[basicstyle=\small\ttfamily,
                   caption={Сигнатура командной строки универсальной версии (\texttt{src/windowed\_dp\_universal.cpp}, репозиторий~\cite{ourgithub})},
                   label={lst:universal-cli}]
./windowed_dp_universal  <SIGMA>  <N>  <MAX_SH>  <K>
\end{lstlisting}
и выполняет следующие шаги:
\begin{enumerate}
    \item Вычисляет $\mathrm{BITS} = \lceil \log_2 \sigma \rceil$, $\mathrm{MASK\_BITS} = \mathrm{BITS} \cdot N$, $M_{\max} = 2^{\mathrm{MASK\_BITS}}$, $M = \sigma^N$ (число валидных масок).
    \item Строит список валидных масок: все целые от $0$ до $M_{\max} - 1$, в которых каждый блок из $\mathrm{BITS}$ бит кодирует число $< \sigma$.
    \item Выделяет динамически массивы слоёв $f_K[\mathrm{sh}][a][b]$ как одномерные \texttt{std::vector<float>} размера $\mathrm{MAX\_SH} \cdot M_{\max}^2$, с обращением только по валидным маскам.
    \item Запускает ДП-итерации по схеме~\S4.5 с заменой $\{0,1\}$ на $\Sigma_\sigma = \{0, \ldots, \sigma-1\}$ и делением на $\sigma^2$ вместо $4$.
    \item По завершении выводит $f_K(0,0,0) / K$ --- нижнюю оценку $\gamma_\sigma$.
\end{enumerate}

Ограничение реализации: $\mathrm{MASK\_BITS} \leq 24$, чтобы $M_{\max}^2 = 2^{2\,\mathrm{MASK\_BITS}}$ помещалось в адекватный объём памяти.

\subsection{Проверка эквивалентности версий}

Универсальная версия при $\sigma = 2$ воспроизводит результаты специализированной (x6) с точностью до округления float (различия в $7$~знаке). Это служит регрессионным тестом для рефакторингов в универсальной версии.

\subsection{Программная инфраструктура}

Все перечисленные ниже файлы расположены в открытом git-репозитории работы~\cite{ourgithub} (\url{https://github.com/Youka4/HSEDiploma}); пути указаны относительно корня репозитория.

\begin{itemize}
    \item \texttt{src/windowed\_dp.cpp} --- специализированная (бинарная) версия оконного ДП, использованная для рекордного запуска при $\sigma = 2$.
    \item \texttt{src/windowed\_dp\_universal.cpp} --- универсальная версия для произвольного $\sigma \geq 2$; принимает параметры $(\sigma, N, \mathrm{MAX\_SH}, K)$ из командной строки (см.\ листинг~\ref{lst:appb-run-commands} в приложении~Б).
    \item \texttt{src/monte\_carlo\_universal.cpp} --- независимая Monte-Carlo верификация классическим ДП-LCS.
    \item \texttt{results/results.txt}, \texttt{results/results\_universal.txt} --- логи запусков с параметрами и полученными значениями $\gamma_\sigma$.
\end{itemize}

Команды запуска серии экспериментов для всех $\sigma \in \{2, \ldots, 10\}$ и команды компиляции бинарников приведены в приложении~Б (листинги~\ref{lst:appb-run-commands} и~\ref{lst:appb-build-commands}).

\subsection*{Выводы по разделу}
\addcontentsline{toc}{subsection}{Выводы по разделу 4}

\begin{enumerate}
    \item Семь итераций оптимизации (x0--x6) позволили перейти от $N = 7$ к $N = 12$ при $\mathrm{MAX\_SH} = 30$.
    \item Ключевые оптимизации: разделение параметров $N$ и $\mathrm{MAX\_SH}$ (x3), переход на одинарную точность и включение AVX2 (x5), развёртка и кэш-локальность.
    \item Универсальная версия (\texttt{src/windowed\_dp\_universal.cpp} в git-репозитории~\cite{ourgithub}) единообразно обрабатывает произвольный $\sigma \geq 2$ за счёт предвычисленного списка валидных масок.
    \item Соответствие версий проверено: универсальная при $\sigma = 2$ воспроизводит специализированную с точностью до округления.
\end{enumerate}
